\documentclass[a4paper]{report}
\usepackage{amsmath} %pour tous les trucs de math
\usepackage{systeme} %pour tous les systemes d'équations
%\usepackage{ bbold }  %pour toutes les doubles lettres
\usepackage{amssymb}  %pour les double Lettres
\usepackage{IEEEtrantools} %pour les équations en collonnes
\usepackage{amsthm} %pour les preuves
\usepackage[english]{babel} % la langue utilisée
\usepackage[utf8]{inputenc} % encodage symboles entrée
\usepackage[T1]{fontenc} % encodage symboles sortie
\usepackage{fancyhdr} %pour les entêtes et pied de page
%\usepackage[math]{blindtext} % pour le Lorem ipsum
%\usepackage{enumitem} %pour changer les listes
\usepackage[a4paper,textwidth=15.5cm,textheight=22cm]{geometry}
\usepackage{graphicx} % pour les graphiques
%\usepackage{subfig} % pour les doubles figures
\usepackage{float} % pour bien positionner les figures
\usepackage[dvipsnames]{xcolor} % pour la couleur du texte et de la page
\usepackage{comment} %pour les blocs de commentaires
\usepackage{biblatex} % bibliographie
\usepackage{csquotes} % pour que la biblio s'adapte à la langue
\usepackage{prettyref}
\usepackage{hyperref} %pour les hyperliens et références(mettre en dernier)
\hypersetup{   
    colorlinks=true,
    linkcolor=blue,
    filecolor=blue,      
    urlcolor=blue,
    citebordercolor= 0 0 0 ,
    linkcolor=blue ,
    citecolor = blue }
\urlstyle{same}
\newrefformat{fig}{Figure~[\ref{#1}]}
\newrefformat{it}{question~\ref{#1}.}
\newrefformat{eq}{(\ref{#1})}
\newrefformat{seq}{Section~\ref{#1}}
\newrefformat{th}{Theorem~\ref{#1}}
\newrefformat{lem}{Lemma~\ref{#1}}
\newrefformat{cor}{Corollary~\ref{#1}}
\newrefformat{rem}{Remark~\ref{#1}}

\newtheorem{theoreme}{Theorem} %[section]
\newtheorem{corollaire}{Corollary} %[theorem]
\newtheorem{lemme}{Lemma} %[theorem]
\theoremstyle{definition}
    \newtheorem{definition}{Definition} %[section]
    \newtheorem*{exemple}{Example}
\theoremstyle{remark}
     \newtheorem*{remarque}{Remark}

\renewcommand{\a}{\alpha}
\renewcommand{\b}{\beta}
\newcommand{\g}{\gamma}
\renewcommand{\d}{\delta}
\newcommand{\m}{\mu}
\newcommand{\n}{\nu}
\renewcommand{\P}{\Phi}
\renewcommand{\l}{\lambda}
\newcommand{\s}{\sigma}
\newcommand{\e}{\epsilon}
\newcommand{\dotx}{\dot{x}}
\newcommand{\doty}{\dot{y}}
\newcommand{\bdotx}{\dot{\mathbf{x}}}
\newcommand{\dotxx}{\dot{\mathbf{x}}}
\newcommand{\dotyy}{\dot{\mathbf{y}}}
\newcommand{\dotzz}{\dot{\mathbf{z}}}
\newcommand{\vv}{\mathbf{v}}
\newcommand{\xx}{\mathbf{x}}
\newcommand{\yy}{\mathbf{y}}
\newcommand{\zz}{\mathbf{z}}
\newcommand{\FF}{\mathbf{F}}
\newcommand{\ww}{\mathbf{w}}
\newcommand{\dotB}{\dot{B}}
\newcommand{\dotP}{\dot{\P}}
\newcommand{\dotI}{\dot{I}}
\newcommand{\dotV}{\dot{V}}
\newcommand{\dotR}{\dot{R}}
\newcommand{\Bstar}{B_{**}}
\newcommand{\Pstar}{\P_{**}}
\newcommand{\R}{\mathbb{R}}
\newcommand{\N}{\mathbb{N}}
\newcommand{\E}{\mathbb{E}}
\newcommand{\C}{\mathbb{C}}
\newcommand{\dd}{\mathrm{d}}
\newcommand{\ddt}{\frac{\dd}{\dd t}}
\newcommand{\ie}{\emph{i.e.\ }}
\newcommand{\ps}[2]{\langle#1,#2\rangle}
\newcommand{\LV}{Lotka-Volterra }

\DeclareMathOperator{\Span}{Span}

\newcommand{\com}[1]{\textcolor{ForestGreen}{[~\emph{#1}~]}}



\pagestyle{fancy}
\lhead{Chapter \thechapter}
\rhead{Section \thesection}
\chead{}

\title{ \includegraphics[scale=0.19]{images/EPFL-logo.png}  \linebreak \linebreak
\bf{Stability of linear dynamical systems and application to Lotka-Volterra equations and epidemiology}}
\author{Benoît Müller \\ \\
Bachelor project in mathematics \\ 
École Polytechnique Fédérale de Lausanne \\ \\
Supervisor: Joachim Krieger \\
Co-supervisor: Shengquan Xiang }
\date{Spring 2021}

\addbibresource{biblio.bib}

\begin{document}
    \maketitle
    \begin{abstract}
  The goal of this project is to study the stability of some linear and nonlinear autonomous ordinary differential systems including the "predator–prey" based Lotka-Volterra models, and to further find applications to epidemiology such as the Covid-19 evolution.
  
   The first chapter is devoted to the investigation of the stability of linear autonomous ODE systems, where a complete classification is given thanks to spectrum information. We refine these results to give a stability decomposition of the solution space, and show that in general, some stability properties are preserved to related nonlinear systems thanks to linearization arguments. 
   
   Then, we turn in the second chapter to a more specific form of ODE systems, say Lotka-Volterra models. Instead of considering classical Lotka-Volterra systems, we look at a more general form. Because at some special critical points the linearization stability theory is not always applicable, we introduced the Lyapunov approach, which can even lead to global stability results. 
   
   Finally, the third chapter is an application of the preceding chapters, where we try to model Covid-19 systems concerning patients and vaccinations, and to predict the evolution of the pandemic via stability investigation.
\end{abstract}

%We introduce linear dynamical systems and investigate the general form of the solutions, as well as the space they form. We introduce the stability of a fixed point, which describe the boundedness or the asymptotic comportment of the distance between a fixed point and a solution starting near it. From the understanding of the solutions, we deduce necessary and sufficient conditions on the matrix that define the equation, to have stability. All conditions relate to the sign of the real part of its eigenvalues. We present a classification of linear dynamical systems, that describe the kind of stability they have. We show some links between linear and nonlinear dynamical systems, such as the linearization of a system. We motivate the study of a particular nonlinear dynamical system on the plane, by showing how it is linked to population dynamics, and present a modification on it that change the affect between the species. We investigate the nature of the solutions and the stability of a particular fixed point. From the study of this problem, and the basic strategic idea of this modelisation, we present adaptations to epidemiology dynamics. We derive implicit formulas for the trajectories of their solutions and some descriptions of the asymptotic comportment.
    \tableofcontents
    \chapter*{Introduction}
\addcontentsline{toc}{chapter}{Introduction} \markboth{INTRODUCTION}{}
In the following paragraphs, we detail the composition of this project and develop the ideas of the main results.
\\\\
In the first chapter we introduce linear dynamical systems, these are autonomous differential systems, where the derivative of a solution is linear with respect to the solution itself. We write it as a matrix product. Solutions always exist on the all real space and are unique for a same initialisation.

We describe the space of solutions. It has the same dimension as the image space because the linear dependence of solutions is the same as their linear dependence at any fixed time. We can compute the solution starting at any point by multiply the initial position by a matrix called fundamental matrix. Its value can be computed as an infinite sum. We allow our self complex solutions, seeing that their real and imaginary parts are real solutions too, and we find the general form of the solutions, by treating particular solutions that start on the generalised eigenvectors of the system matrix. They are called eigensolutions and are the product of two factors: an exponential and a polynomial. The eigenvalue appears in the exponential and the polynomial is of coefficients generalised eigenvectors of the same eigenspace. A basis of generalised eigenvectors form a basis of eigensolutions for the solution space.

Then we introduce the stability of a fixed point. We say it is Lyapunov-stable when all solutions that start near enough the fixed point never leave an arbitrary neighbourhood. It is asymptotically stable if in addition to that, all solutions that start near enough the fixed point, converge to it. From the understanding of the solutions, we deduce necessary and sufficient conditions on the matrix that define the equation, to have stability on the origin. All conditions relate to the sign of the real part of its eigenvalues, which appear in the argument of the exponential of the eigensolutions. It is asymptotically stable if and only if all eigenvalues have a negative part, and we call such eigenvalues stable. It is Lyapunov-stable if and only if all eigenvalues have a non-positive real part, and a negative one when the eigenvalue is defective. We give a classification of linear dynamical systems, that describe the kind of stability they have. They take into account the range of stability of the eigenvalues.

We present a direct sum that decompose the space of solutions in three subspaces: the generalised eigenspaces associated to eigenvalues of real parts that are respectively negative, null, and positive. These three spaces, called respectively stable, center and unstable eigenspace have a second description that only take the asymptotic comportment of solutions when t tends to plus or minus infinity. Two system that have a trivial center eigenspace are actually linked by a homeomorphism between their solutions. We show a direct application of our understanding of the linear case, the linearization theorem. It says that if the linearization of the system in a fixed point is asymptotically stable, then the nonlinear system is asymptotically stable too. 
\\ \\
We motivate the study of a particular nonlinear dynamical system on the plane, called \LV equations, by showing how it is linked to population dynamics. The functions that are set in an equation are the size of two population of prey and predators. It is based on the concept of logistic growth and the understanding of the growth rate, which is the derivative divided by the function itself. It represent the mean of change an individual add to the population by unit of time. This rate of growth is linear in the \LV equation. We present a modification on it, that change the affect between the species. The idea is to consider in the equation not just the effective size of population, but a function of it. This is useful to have a better description of the affect each species have on the others. We consider two situations. In the first system, a function does not appears in its own growth rate, and in the second it does. 

We investigate the nature of the solutions, for the first system, the linearization in a particular fixed point gives an elliptic linear system, a system with cyclic solutions. We search then for a constant quantity that is conserved along time, and this gives a implicit equation for the trajectory of the solutions. From this, we conclude that the solutions are indeed cyclic, they orbit around the fixed point. In the second system, the linearization in a particular fixed point gives an stable linear system, so we know that the fixed point is asymptotically stable. Motivated by this result, and inspired by the conserved quantity of the first system, we develop the theory of Lyapunov functions. They are function that represent the energy of a position, and such that all solutions loose energy along time. When such a function exist and has a minimum in a fixed point, we can obtain asymptotic stability without condition on the starting point of the solutions, they all converge to the fixed point. Taking a similar expression of the constant quantity of the first system, we find a Lyapunov function. From this, we conclude that all the solutions turn in a spiral around the fixed point and converge to it, we have global asymptotic stability.
\\ \\
From the study of this problem, and the basic strategic idea of this modelisation, we present adaptations to epidemiology dynamics. We investigate the modelisation of a population that is confronted to a virus. We use the \LV system to describe the population of vulnerable people that have not been infected, the population of infected people, and the population of people that have recovered after infection. We see that this model has no isolated fixed point so fixed points are not asymptotically stable. We derive again constant quantities along time, that give implicit formulas for the trajectories of solutions and the limits of them. If the initial number of vulnerable people is bigger than a certain threshold, the number of infected people increase to a maximum located in this threshold, and eventually tends to zero. If the initial number of vulnerable people is smaller than this threshold, the number of infected people directly tends to zero. We give some description of the asymptotic comportment and the reminded size of groups.
In a second modelisation, we introduce the vaccine to the population and separate each group in two, regarding if they are vaccinated or not. The analysis of the problem can be done with the same tools as the precedent, with a bit more work. We obtain a pair of equations that gives some explicit relations between the functions. In particular, we have a bijective relation between the vulnerable groups and show the affect of the vaccine.
\thispagestyle{plain}

    \chapter{Linear dynamic systems and stability}
In this chapter, we study explicit autonomous ordinary differential equations of first order that are linear, namely equations of the form 
\[\dotxx=A\xx\]
where $A$ is a $n\times n$ matrix and $\xx$ is a differentiable function from a real interval to $\R^n$. since its derivative should be $C^1$ too, $\xx$ must be $C^\infty$ to exist.

Our main goal is to investigate stability. For this, we present a brief description of the solution space and explicit forms of its elements. This first section will help us for the second, in the study of the stability.
\section{Solution space}
First of all, using results by Cauchy and Picard, and the particular smoothness of $\xx\to A\xx$, we know that the solutions locally exist, and that they are unique for a same initialisation. Since $$\ddt\|\xx\|^2=2\ps{\xx}{\dotxx}=2\ps{\xx}{A\xx}\leq 2\|A\|\|\xx\|^2,$$ 
$$\|\xx(t)\|^2\leq \|\xx(0)\|^2e^{2\|A\|t} $$
by Grönwall and solutions cannot explode, they are defined on all $\R$. We note that since the derivative of a solution is just itself multiplied by a matrix, it is derivable too, and by induction $\xx$ is $C^\infty$ with $\xx^{(k)}=A^k\xx$. 

The equation is autonomous, so any solution can be reparametrised by a translation of the time:
If $\dotxx(t)=A\xx(t)$, then $\ddt(\xx(t+\tau))=A\xx(t+\tau)$. That is why we will generally initialise them at $t=0$ without loss of generality. Solutions are described by a \emph{flow} $\phi:\R^n\times \R\to\R^n$ where $t\mapsto\phi(\xx_0,t)$ is the solution starting in $\xx_0$ at $t=0$. However, existence theorems does not tell us much about their actual form. We present then their existence in this particular case.
First, we observe that the set of solution is a vector subspace of $C^\infty$. Indeed for all solutions $\xx$, $\yy$ and a scalar $\a$, \[(\a\xx+\yy)'=\a\dotxx + \dot{\yy}=\a A\xx + A\yy=A(\a\xx+\yy)\] and the identical null function is trivially in the space. This motivate us to find a basis of this subspace and understand how to construct it.
\begin{definition}
    A collection of $k$ solutions $\xx_1,\dots,\xx_k$ on $I$ are said to be \emph{linearly independent} or \emph{independent}, if $\xx_1(t),\dots,\xx_k(t)$ are linearly independents for all $t\in I$.
\end{definition}
\begin{lemme} \label{lem:indépendance}
For any scalar $\tau$, a collection $\xx_1,\dots,\xx_k$ of solutions are linearly independents if and only if $\xx_1(\tau),\dots,\xx_k(\tau)$ are linearly independents.
\end{lemme}
\begin{proof}
If the solutions are independent we have the result by definition. In the other way, by uniqueness of the solutions up to initialisation, we have that if for time $\tau$, the positions are not independent, then we must have non all null scalars $\a_1,\dots,\a_k$ such that  $\a_1\xx_1(\tau) + \dots + \a_k\xx_k(\tau)=0$. But then 0 is a trivial solution starting there and actually  $\a_1\xx_1(t) + \dots + \a_k\xx_k(t)=0$ for all $t$.
\end{proof}
This lemma prove that the space of solution has the same dimension as the space $\R^n$, \ie $n$. 
In order to manipulate the solutions, we put them in a matrix, like $(\xx_1,\dots,\xx_k)$. In this form we see that we can actually extend the equation to matrix entries: 
\[\dot{X}=AX\]
and since 
\[\dot{X}=(\dotxx_1,\dots,\dotxx_k)  
\text{ and } 
AX=(A\xx_1,\dots,A\xx_k),\]
a matrix X is a solution if and only if its columns are vector solutions. We see that for a vector $v\in\R^k$, $Xv$ is a vector solution $(Xv)' = \dot{X}v = AXv$. This solution ca be written as $Xv = v_1\xx_1 + \dots + v_k\xx_k$ and is then a linear combination of the column solutions of $X$. Thus, a matrix solution allow us to easily write new solutions as linear combinaison of a collection of solutions. As the dimension of the space of solutions is $n$, exactly $n$ linear independents solutions will be enough to construct all solutions as a product of the matrix solution constructed with them.
\begin{definition}
    If a matrix solution $M=(\xx_1,\dots,\xx_n)$ is square (k=n) and of full rank, then it is called a \emph{fundamental matrix solution} and $\{\xx_1,\dots,\xx_k\}$ a \emph{fundamental set of solutions}.
\end{definition}
With such a $M$, $M\R^n$ is the all space of solutions. At time $t=0$, $M(0)v=v_1\xx_1(0) + \dots + v_k\xx_k(0)$ meaning that with the condition $\xx(0)=x_0$, v must be chosen to be the vector of coordinates of $x_0$ in the basis of the columns of M.
\\ \\
 We can now look at the forms the solutions can take. In one dimension, the problem is $\dotx=ax$ and if a is non null, non trivial solutions satisfy $a=(\log|x|)'$ which give $\log|x(t)|=at+c$ and $x(t)=\pm e^{at+c}=x_0e^{at}$. We check easily that $e^{at}$ is indeed a solution. The $n$-dimensional case is more complicated. We try to put the equation in its integral form and derive a property by recurrence (supposing $X_0=I$ for simplicity): 
\begin{IEEEeqnarray*}{rCl}
X(t) 
&=& I + \int_0^t AX(s) \dd s 
= I + \int_0^t A\bigg(I + \int_0^{s_1} AX(s_2) \dd s_2\bigg) \dd s_1
= I + tA +\int_0^t \int_0^{s_1} A^2X(s_2) \dd s_2 \dd s_1
\\ &=& I + tA + \frac12(tA)^2 + \int_0^t \int_0^{s_1} \int_0^{s_3} A^3X(s_3) \dd s_2 \dd s_1.
\end{IEEEeqnarray*}
We see the Taylor expansion of $\exp$ appear, as a generalisation to matrices, where we define $A^0=I$ for any matrix $A$. This motivate the following definition:
\begin{definition}
If it converges, we define the \emph{exponential} $e^A$ of a square matrix A and the underlying function $\exp$ as the infinite sum 
\[e^A = \sum_{n=0}^\infty \frac{1}{n!}A^n.\]
\end{definition}
\begin{lemme} \label{lem:exp}
The exponential of a matrix always converge.
\end{lemme}
\begin{proof}
We use $\|.\|$ for the operator norm on matrices, keeping in mind that all matricies norms are equivalent. Now we have by basic properties of this norm that 
\[ \sum_{n=0}^N \|\frac{1}{n!}A^n\| 
\leq \sum_{n=0}^N \frac{1}{n!}\|A^n\| 
\leq \sum_{n=0}^N \frac{1}{n!}\|A\|^n ,\]
and this is the Taylor finite expansion of $\|A\|$, which converge for all value of $\|A\|$. As a result the sum is absolutely convergent for the operator norm, and then must converge.
\end{proof}
\begin{theoreme}
For a matrix $A$ and a scalar $t$, the quantity $e^{At}$ is differentiable with respect to $t$ and has derivative $\dd/{\dd t} (e^{At}) = Ae^{At}$.
\end{theoreme}
\begin{proof}
Each coordinate of the series is actually a convergent Taylor series in t and is then analytic. The theory of analytic functions tell us that they are $C^\infty$ and that we can derive term by term.
%However, we explain the argument of the proof in our case: As seen in \prettyref{lem:exp}, the sum converge absolutely, and in our case it converge locally uniformly in $t$ when $N\to\infty$, because if $|t|\leq t_{max}$, 
%\[ \sum_{n=0}^N \|\frac{1}{n!}t^nA^n\|
%\leq \sum_{n=0}^N \frac{1}{n!}(t_{max}\|A\|)^n\] converge without independence in $t$ for the last quantity. 
 We get 
 \[ \ddt e^{At} =\frac{\dd}{\dd t}\sum_{n=0}^\infty\frac{1}{n!}t^nA^n
= \sum_{n=1}^\infty \frac{1}{(n-1)!}t^{n-1}A^n
= A\sum_{n=0}^\infty \frac{1}{n!}(tA)^n
=Ae^{At} \]
%converge locally uniformly to $Ae^{tA}$ for the same reason. As a result, since the partial sums and the derivatives of the partial sums converge locally uniformly, the limit is differentiable with derivative the limit of the derivatives of the partial sums, namely $Ae^{tA}$ and we have the result.
\end{proof}
All of this tell us that as we expected, $e^{tA}$ is a matrix solution to the linear differential equation.
\begin{corollaire}
The matrix $e^{tA}$ is a fundamental solution and each solution write $\xx(t)=e^{tA}\xx(0)$
\end{corollaire}
\begin{proof}
We evaluate it in $t=0$ by direct calculation since the sum become finite: 
\[e^{0A}=\sum_{n=0}^\infty \frac{1}{n!}(0A)^n = I.\]
So the initial matrix is non singular, meaning that it will stay along the time non singular and $e^{tA}$ is a fundamental matrix solution with formula $x(t)=e^{tA}x_0$.
\end{proof}

For matrices that are not in a special form, we cannot directly compute the series of the exponential and then then neither the solutions. Specials forms of matrices whose exponential are commutable are for example the ones where the power of the matrix has a general formula, like diagonal matrices. We know that diagonal matrices are directly related to eigenvalues and eigenvectors. The eigenvector are the directions where the matrix act like in the one dimensional case and as often, it is a way to find similar results in the higher dimensional problems.

Let us investigate what happens in theses directions, and choose a candidate that look like a one dimensional solution,$e^{\l t}v$, for a real eigenvalue $\l$ of the matrix $A$ and the corresponding eigenvector~$\vv$. We compute its derivative and obtain 
\[\dd/{\dd t}(e^{\l t}v) =  e^{\l t}\l v = e^{\l t}Av = A(e^{\l t}v).\]
This give us indeed a non trivial solution since the eigenvector is non null. Theses kind of solutions can be seen in the computation of the exponential when we have a basis of eigenvectors: the matrix $A$ is diagonalisable like $A=PDP^{-1}$ where $D$ is the diagonal matrix of the eigenvalues and P is the non singular matrix formed by the eigenvectors. Then 
\[ A^n= (PDP^{-1})^n = PDP^{-1}\cdots PDP^{-1} = PD^nP^{-1} \]
and 
\[ e^{tA} 
= Pe^{tD}P^{-1} 
= P\text{diag}(e^{\l_jt})P^{-1} 
= (v_1e^{\l_1t},\dots,v_ne^{\l_nt})P^{-1} 
\]
which is the fundamental matrix of the kind of solutions we find before, up to a reparametrization.

This was the simple case. More generally the matrix can have complex eigenvalues with complex eigenvectors, but the following lemma show us that complex solutions are made of real solutions :
\begin{lemme} \label{lem:complex}
If $\zz=\xx+i\yy$ is a complex solution of real and imaginary parts $\xx$ and $\yy$, then $\xx$ and $\yy$ are real solutions.
\end{lemme}
\begin{proof}
This can be shown by the simple fact that
\[\dotxx+i\dotyy =\dotzz = A\zz = A\xx+iA\yy. \]
As A is real, we can directly read the real and imaginary parts as being $A\xx$ and $A\yy$. The real and respectively imaginary parts being the same along equalities implies $\dotxx=A\xx$ and $\dotyy=A\yy$ which is the result.
\end{proof}
We must then consider complex solutions obtained by complex eigenvalues too. We sum up the results together in the following theorem.
\begin{theoreme} \label{th:eigensolutions}
For a complex eigenvalue $\s=\a+i\b$ of $A$ with eigenvector $w=u+iv$, we have the real solutions
\[e^{\a t}(\cos(\b t)u - \sin(\b t)v)\] 
\[e^{\a t}(\cos(\b t)v + \sin(\b t)u).\] 
They are independent if ad only if $u$ and $v$ are independent. If $\s$ is real, this give only one solution $e^{\a t}u.$
\end{theoreme}
\begin{proof}
First we have that $e^{\s t}w$ is a complex solution, since 
\[\ddt(e^{\s t}w) = e^{\s t}\s w = e^{\s t}Aw = A(e^{\s t}w).\]
Then by \prettyref{lem:complex}, the real and the imaginary parts are real solutions, we compute them by rewriting the complex solution:
\[e^{\s t}w = e^{\a t}(\cos(\b t) + i\sin(\b t))(u+iv)
= e^{\a t}(\cos(\b t)u - \sin(\b t)v) + ie^{\a t}(\cos(\b t)v + \sin(\b t)u).\]
We recognise the resulting solutions in the real and imaginary part. If we evaluate them in t=0, we get respectively $u$ and $v$ so the independence follow by \prettyref{lem:indépendance}. By setting $\b=0$, the solutions are $e^{\a t}(\cos(0)u - \sin(0)v) = e^{\a t}u$ and  $e^{\a t}v$. But since the eigenvalue is real, the eigenvector is real too and $v=0$ letting only one non-trivial solution.
\end{proof}
\begin{remarque}
Note that complex values come by pairs of conjugates, as well as the eigenvectors:
$Aw=\s w$ implies 
$$A\overline{w} = \overline{A}\overline{w} = \overline{Aw}= \overline{\s w} = \bar{\s}\bar{w}.$$
However, this will not give us new real solutions by the method of \prettyref{th:eigensolutions} because $$e^{\overline{\s} t}\overline{w}= \overline{e^{\s t}}\overline{w} = \overline{e^{\s t}w}$$
has the same real and imaginary part as $e^{\s t}w$ up to the sign. As a result, a complex eigenvalue give two solutions but together with its conjugate, they give us still two solutions, so we can find as many independent solutions as independent eigenvector we find, real or not.
\end{remarque}
Now we have to deal with the case when we do not have a basis of eigenvectors. In this case, some eigenvalue $\l$ with algebraic multiplicity $\m_a$ and a geometric multiplicity
$$\m_g=\dim\ker(A-\l I) < \m_a.$$
Such an eigenvalue is called \emph{defective}, and the matrix is said \emph{defective} too when it has at least one defective eigenvalue. We use here the concept of generalized eigenvector that come from the result about the Jordan form:

\begin{definition}
    A vector $\ww$ is a \emph{generalized eigenvector} of rank $m$  of a matrix A and corresponding to an eigenvalue $\l$, if it is a vector (possibly complex if $\l$ is complex) that satisfy
    $$(A-\l I)^m\ww=0 \quad\text{and}\quad (A-\l I)^{m-1}\ww\neq0$$
    for $m\in\N^*$.
    
    A \emph{canonical} basis of generalised eigenvectors is a basis of generalised eigenvector such that for all generalized eigenvector $\ww$ of rank $m$ that is in the basis, for all $0<j<m$, $(A-\l I)^j\ww$ are generalised eigenvectors of rank $m-j$ with respect to $\l$, and are in the basis too.
\end{definition}

\begin{remarque}
As before, a square matrix powered by 0 is the identity matrix. Note that $(A-\l I)^{m-1}\ww\neq0$ implies that in particular $\ww\neq0$. As a consequence, a generalised eigenvector of rank 1 is a usual eigenvector.\com{crossed out in the review but have not seen why}
\end{remarque}
    
\begin{theoreme}
    A matrix always have a canonical basis of generalised eigenvectors.
\end{theoreme}
\begin{proof}
    Linear algebra result, see \com{reference} for a proof.
\end{proof}

In term of computation, we start from $(A-\l I)\vv=0$ to find the eigenvectors. Then we search for some $\ww$ such that $(A-\l I)\ww=\vv$, assuring that $(A-\l I)^2\ww=(A-\l I)\vv=0$ and hence that $\ww$ is a generalised vector of rank 2. We continue like this making a chain of generalised vectors. Even if we decide to take into account complex vectors

We show that this basis is useful, if we use its components as initialisation points.
\begin{theoreme}[General form of the eigensolutions] \label{th:solutiondegeneree}
 For a canonical basis $B$ of generalised eigenvectors, we have a basis of complex solutions with form $e^{t\l}p_\ww(B)$, where $\l$ is the eigenvalue associated to a generalised eigenvector $\ww$ of rank $m$ in the basis, and $p_\ww$ is a polynomial of degree $m-1$ with coefficients that are in the eigenspace of $w$ .
\end{theoreme}
\begin{proof}
We see easily that the basic property of the real exponential that changes sum into product is true for the matrices with the condition that they commute. The proof is the same as in the real case as we have already proven absolute convergence. 
We write $A= \l I + (A-\l I) $ and $\l I$ is diagonal hence commutative with all matrices. Taking $\ww$ from a canonical basis, we get
\begin{IEEEeqnarray}{rCl} \label{eq:solutioncomplexe}
e^{tA}\ww 
&=& e^{t\l I + t(A-\l I)}\ww 
=e^{t\l I} e^{t(A-\l I)}\ww 
= e^{t\l}\sum_{n=0}^\infty \frac{1}{n!}t^n(A-\l I)^n\ww
= e^{t\l}\sum_{n=0}^{m-1} \frac{1}{n!}t^n\ww_n 
\end{IEEEeqnarray}
where the $\ww_n=(A-\l I)^n\ww$ are other generalised vectors by definition of the canonical basis. So have indeed a polynomial in $t$ with generalised eigenvectors coefficients. These are surely independent solutions for different choices of $\ww$ since the generalised eigenvectors are supposed independent and they are the initial values of these solutions, the proof is complete.
\end{proof}
We have now a complex basis but we know that eigenvalues and eigenvectors come by pairs of conjugates. It is the same for generalised eigenvectors and eigenvalues:
$$(A-\l I)^m\ww=0 \quad\text{and}\quad (A-\l I)^{m-1}\ww\neq0$$
implies
$$(A-\overline{\l} I)^m\overline{\ww}=0 \quad\text{and}\quad (A-\overline{\l} I)^{m-1}\overline{\ww}\neq0$$
So we will get solutions like \prettyref{eq:solutioncomplexe} by pairs of conjugates $e^{tA}\ww$ and $e^{tA}\overline{\ww}$. by subtracting them or summing them, we get two real new solutions that we will call \emph{degenerated}:
$$
    e^{tA}\ww + e^{tA}\overline{\ww} 
    = e^{tA}(\ww+\overline{\ww})
    = 2e^{tA}\Re(\ww), 
$$
$$
    e^{tA}\ww - e^{tA}\overline{\ww} 
    = e^{tA}(\ww-\overline{\ww})
    = 2e^{tA}\Im(\ww).
$$
However, we can see that the formulas for the real and the imaginary part are not very convenient:
\begin{IEEEeqnarray*}{C}
e^{t\l}\sum_{n=0}^{m-1} \frac{1}{n!}t^n\ww_n \\
= e^{t\a}(\cos(\b t)+i\sin(\b t))\sum_{n=0}^{m-1} \frac{1}{n!}t^n(\Re\ww_n + i\Im\ww_n) \\
= e^{t\a}\Big(\cos(\b t)\sum_{n=0}^{m-1} \frac{1}{n!}t^n\Re\ww_n-\sin(\b t)\sum_{n=0}^{m-1} \frac{1}{n!}t^n\Im\ww_n\Big)
+ ie^{t\a}\Big(\cos(\b t)\sum_{n=0}^{m-1} \frac{1}{n!}t^n\Im\ww_n+\sin(\b t)\sum_{n=0}^{m-1} \frac{1}{n!}t^n\Re\ww_n\Big).
\end{IEEEeqnarray*}
We know that we have a (complex) basis of the complex space solution. The dimension of the $\C$-vector space $\C^n$ is $n$, and since $\R\subset\C$, all real solutions are included in this description. However, the method of taking the real part of a basis does not give us a priori a independent set of real solutions. For this reason, we will stay in the the complex space because it present a good and simple description thanks to its algebraic closure, and we will be able too deduce from the complex case results about what we need for the real case. We state a corollary of the precedent theorem, giving all possible forms that can take the reals solutions, but without giving a explicit formula.
\begin{corollaire} \label{cor:formesolutionlineaire}
    All solutions to linear system have coordinates that are linear combinations of the following functions :
    \begin{itemize}
    \item $e^{\l t}$ 
    \item $e^{\a t}\cos(\b t)$ , $e^{\a t}\sin(\b t)$
    \item $t^je^{\l t}$ , $t^je^{\a t}\cos(\b t)$ , $t^je^{\a t}\sin(\b t)$ 
    \end{itemize}
    where $\l$ is a real eigenvalue, $\s=\a+i\b$ is an complex eigenvalue, $0\leq j< m_a$ is a natural number with $m_a$ the algebraic multiplicity of $\s$. Note that each point is a generalisation of the precedent.
\end{corollaire}
\begin{remarque}
    This actually describe the complex space too, since we can rebuild the complex solutions we used to derive this corollary by summing the real part and $i$ times imaginary one.
\end{remarque}
\section{Stability}
In other terms, \prettyref{th:eigensolutions} tell us that each non null real eigenvalue gives direction(s) where the trajectories are straight and of exponential velocity. Null eigenvalues give direction(s) where trajectories are fixed. Non real eigenvalue gives direction(s) where the trajectories are like ellipses that can change of size exponentially. The \prettyref{th:solutiondegeneree} add other sorts of solutions as polynomials resized by a exponential and correspond to the case of degenerated eigenvalues.

All these considerations are only on the special directions. Depending on the sign of the real part of the eigenvalues, theses specials solutions go very fast to $0$, stay in an orbit, diverge very fast to infinity values, or maybe follow a polynomial. This motivate us to see how these special solutions act together, and what is the asymptotic comportment of trajectories, how stable is the origin, and make a classification about all theses factors.
First of all we define concepts about stability, and asymptotic convergence. For this we put our self in a more general cadre which will be useful later, with function $\FF$ and the equation 
\[\dotxx = \FF(\xx)\]
We suppose $\FF$ regular enough to have a flow $\phi(\xx_0,t)$ which encapsulate all solutions:
\begin{IEEEeqnarray*}{rCl}
\dot{\phi}(\xx_0,t) &=& \FF(\phi(\xx_0,t)) \\
\phi(\xx_0,0) &=& \xx_0
\end{IEEEeqnarray*}
\begin{definition}
    A solution x is \emph{Lyapunov stable} or \emph{L-stable} if the continuity of the flow with respect to the initial condition $\xx_0$ is uniform with respect to positive times $t\geq$. Namely, if for all $\e>0$ there exist a $\d>0$ such that $\|\zz-\xx(0)\|<\d$ implies that for all $t\geq0$, $\|\phi(\zz,t)-\xx(t)\|<\e$.
\end{definition}
\begin{definition}
    Two solutions $\xx$ and $\yy$ are $\omega$\emph{-attracted} to each other if $\lim_{t\to\infty}\|\yy(t)-\xx(t)\| = 0$. The resulting equivalence class is called the \emph{basis of attraction}. A solution $\xx$ is said \emph{$\omega$-attracting} or \emph{attracting} if there is a $\d>0$ such that $\|\zz_0-\xx_0\|<\d$ implies that $\phi(\zz_0,t)$ is in the same basin of attraction of $\xx$. A solution is said \emph{globally $\omega$-attracting} (or \emph{globally attracting}) on a set if this set is in the basin of attraction. We do not need to specify the set if it is the all solution space.
\end{definition}
\begin{definition}
    A solution is \emph{asymptotically stable} if it is L-stable and attracting. It is said \emph{globally asymptotically stable} if it is L-stable and globally attracting
\end{definition}
\begin{remarque} \label{rem:stabilité}
Note that these two notions are different. For an example of non-attracting point which is L-stable, we can take simply $F(x)=0$ or for a non trivial case, we can take $\dotx = -y$, $\dot{y} = x$ who describe the circle trajectories $x(t)=cos(t)$, $y(t)=sin(t)$.

In the other way there exist non L-stable points which are attracting. Such a point is the limit of all near trajectories but they always go far before converging, like a detour. For this we place our-self in polar coordinates. We want the trajectories to follow the circle and finish in $(1,0)$ for this we make the $\theta$ always go and stop to $2\pi$, and $r$ go and stop to 1. For this we can write $(\dot{r},\dot{\theta})=(1-r,2\pi-\theta) = G(r,\theta)$. But if we want $(G_1\cos(G_2),G_1\sin(G_2))$ to be continuous on $\R_+\times 0$, we should write $\dot{r} = r(1-r)$, $\dot{\theta} = \theta(2\pi-\theta)$. And to obtain the continuity of the derivative, in purpose to have a flow, we could write $\dot{r} = r(1-r^2)$, $\dot{\theta} = \theta(2\pi-\theta^2)$. This gives us what we need but to be able to explicitly change the coordinates into cartesian, we prefer $\dot{\theta} = \sin(\theta/2)^2$. The fact that we started as a polar coordinate and that they are independent actually give a proof that all solutions tends indeed to $(1,0)$ but solutions beginning near $(1,0)$ and in the upper plane, pass behind the origin before returning to the fixed point $(1,0)$. The show that it is attracting but not L-stable. \com{graphic}
\end{remarque}
\begin{lemme}
    When $\FF$ is linear, i.e. $\dotxx=A\xx$, all solutions have the same L-stability and attractivity.
\end{lemme}
\begin{proof}
In the linear case, $\phi$ is linear in the first variable, indeed $\phi(\xx_0,t) = e^{tA}\xx_0 =: X(t)\xx_0$. Now for any solution $\xx$ and all initial point $\zz$,
\[ \|\phi(\zz,t)-\xx(t)\| = \|X(t)\zz-X(t)\xx(0)\| = \|X(t)(\zz-\xx_0)\| =\|\phi((\zz-\xx_0),t)-\phi(0,t)\|, \] meaning that L-stability and attractivity is entirely determined by the stability of the trivial solution $\phi(0,t)=0$.
\end{proof}
\begin{remarque}
    Therefore, we can speak of the L-stability and the attractivity of a linear system, meaning that it applies to all solutions, or none of them, and doing it by looking at the trivial solution 0. In this case the system is L-stable if and only if there exists a $\d>0$ such that $\|\xx_0\|<\d$ implies that for all $t\geq0$, $\|\phi(\xx_0,t)\|<\e$. The system is attracting if and only if there exists a $\d>0$ such that $\|\xx_0\|<\d$ implies that $\phi(\xx_0,t) \to 0$.
\end{remarque}
 
 \begin{theoreme} \label{th:stablecondition}
     The linear system is L-stable if and only if each of its solutions is bounded for positive times.
 \end{theoreme}
 \begin{proof}
  Suppose the system is L-stable, and for contradiction that a solution $\xx$ is not bounded. Let $\d>0$ be the distance given by the stability, such that all solutions that start with a norm smaller than $\d$ do not go away the unit ball. We can then define an other solution $\yy=\d\xx/\|\xx(0)\|$ and $\|\yy(0)\|=\d$, so $\yy$ wont go away the unit ball and is bounded. This contradict the fact that $\xx$ and $\yy$ are proportionals and all solutions must be bounded.
  
  Suppose now that all solutions are bounded. Then the columns of $X=e^{tA}$ the fundamental system are bounded, implying that actually the norm of X is bounded (all norm on finite dimensional spaces are equivalent). Now we get 
  \[ \|\phi(x_0,t)\|=\|X(t)x_0\| 
  \leq \|X(t)\| \|\xx_0\| 
  \leq \max_{t\geq0} \|X(t)\| \|\xx_0\| \]
  which is smaller than any positive $\e$ as soon as 
  $$\|\xx_0\| < \frac{\e}{\max_{t\geq0} \|X(t)\|}$$
  and give us the stability of the system.
 \end{proof}
\begin{theoreme}
    The linear system is globally asymptotically stable if and only if it is attracting.
\end{theoreme}
\begin{proof}
By definition, global asymptotic stability implies global attractivity and so in particular attractivity with any radius condition. In the other direction, if it is attracting with radius condition $\d$, such that when $\|\xx_0\|<\d$, $\|\phi(\xx_0,t)\|\to0$. Then any solution $\xx$ can be written $\xx=\|\xx(0)\|\yy/\d$ where $\yy=\d\xx/\|\xx(0)\|$ is a proportional solution that start with norm $\d\|\xx(0)\|/\|\xx(0)\|=\d$ and is small enough to converge to zero, implying that $\xx=\xx(0)\yy/\d$ converge to zero too. The system is globally attracting by arbitrarity of $\xx$. Now that all solutions are converging to zero, they are all bounded and by \prettyref{th:stablecondition}, we know that the system is actually L-stable. Both condition of stability and global attractivity are reunited, the system is globally asymptotically stable.
\end{proof}

\begin{exemple}[saddle]
Consider the system
$$\begin{cases}
    \dotx=x+3y \\
    \doty=3x+ay
\end{cases}
\text{ rewritten with } \zz=(x,y)^\top: \quad \dotzz=
\begin{pmatrix} 1&3\\3&1 \end{pmatrix}\zz =A\zz
$$
The characteristic polynomial is $p_A(\l)=\det(A-\l I)=(1-\l)^2-9=(-2-\l)(4-\l)$ and has roots $\l_1=-2$, $\l_2=4$. For the eigenvalue -2, $$ A-\l_1I = A+2I= \begin{pmatrix} -3&3\\3&-3 \end{pmatrix} $$
has $\begin{pmatrix}1\\-1\end{pmatrix}$ in the kernel and this is the first eigenvector. The stable associated eigensolution is $$e^{-2t}\begin{pmatrix}1\\-1\end{pmatrix} \xrightarrow[t\to\infty]{} 0.$$
Similary, 
$$ A-\l_2I = A-4I= \begin{pmatrix} 3&3\\3&3 \end{pmatrix} $$
give the eigenvector $\begin{pmatrix}1\\1\end{pmatrix}$ and the unstable eigensolution $e^{4t}\begin{pmatrix}1\\1\end{pmatrix}\xrightarrow[t\to\infty]{\|.\|}~\infty$. We plot the results: \com{graphique + commentaires}
We see that we have only two straight solutions, the eigensolutions. Their norm evolve exponentially but they are time inverted from each other. All the other solutions are combinations of those two and are hyperbolas, they come near the origin and then goes back away. The only one that goes to zero is the eigensolution with eigenvalue -2. We can be as near to zero as we want for the start, the solution eventually goes to infinity if it is not this eigensolution. We does not see any stability, a such system is called a saddle and we will define a general definition for it later.
\end{exemple}

\begin{exemple}[focus]
Consider the system
$$ \dotzz=
\begin{pmatrix} -4&-2\\3&-11 \end{pmatrix}\zz =A\zz
$$
The characteristic polynomial is $(-4-\l)(11-\l)+6=(-5-\l)(-10-\l)$ and has roots $\l_1=-10$, $\l_2=-5$. For the eigenvalue -10, $$ (A-\l_1I)\begin{pmatrix}1\\3\end{pmatrix}
= \begin{pmatrix}6&-2\\3&-1\end{pmatrix} \begin{pmatrix}1\\3\end{pmatrix} = 0.$$
The stable associated eigensolution is $e^{-10t}\begin{pmatrix}1\\3\end{pmatrix} \xrightarrow[t\to\infty]{}~0$. Similary, 
$$ (A-\l_2I)\begin{pmatrix}2\\1\end{pmatrix}
= \begin{pmatrix} 1&-2\\3&-6 \end{pmatrix} \begin{pmatrix}2\\1\end{pmatrix} = 0$$
give the eigenvector $\begin{pmatrix}1\\2\end{pmatrix}$ and the second stable eigensolution $e^{-5t}\begin{pmatrix}1\\2\end{pmatrix}\xrightarrow[t\to\infty]~0$. We plot the results: \com{graphique + commentaires}
Here the two eigensolutions goes to zero exponentially and thus, all linear combinations do the same. The system is attracting and so asymptotically stable. We will call this situation later a focus, and more generally a stable node.
\end{exemple}

\begin{exemple}[null eigenvalue]
Consider the system
$$ \dotzz=
\begin{pmatrix} -1&-3\\-1&-3 \end{pmatrix}\zz =A\zz
$$
The characteristic polynomial is $(-1-\l)(-3-\l)-3=\l(-4-\l)$ and has roots $\l_1=-4$, $\l_2=0$. The eigenvectors are respectively $\begin{pmatrix}1\\1\end{pmatrix}$ and 
$\begin{pmatrix}3\\-1\end{pmatrix}$. They give the solutions $e^{-4t}\begin{pmatrix}1\\1\end{pmatrix}$ and $\begin{pmatrix}3\\-1\end{pmatrix}$ (constant solution). \com{graphique + commentaires}
We see that we have a line where the solutions are fixed and a other where the solution goes exponentially to the origin. As a result the system is L-stable because solutions do not go away, but it is not attracting or asymptotically stable because this fixed line contain solutions that never move to the origin.
We plot the results: \com{graphique + commentaires}
\end{exemple}
These are simple examples where the matrix is diagonalisable. When it is not the case, the situation is more complicated because of the generalised eigenvectors:
\begin{exemple}[defective eigenvalue]
Consider the system
$$ \dotzz=
\begin{pmatrix} -2&1\\-1&0 \end{pmatrix}\zz =A\zz.
$$
The characteristic polynomial is $(1+\l)^2$ and has one root $\l=-1$, $\l_2=0$; $(A+I)=\begin{pmatrix} -1&1\\-1&1 \end{pmatrix}$ give the eigenvector $\vv=\begin{pmatrix}1\\1\end{pmatrix}$. It is the only one because $\det \begin{pmatrix} -1&1\\-1&1 \end{pmatrix} =0$ and it must be of rank one. This give us a solutions $e^{-t}\begin{pmatrix}1\\1\end{pmatrix}$. We search for a generalised eigenvector:
$$(A+I)\ww=\vv \text{ , \ie } \begin{pmatrix} -1&1\\-1&1 \end{pmatrix}\ww= \begin{pmatrix}1\\1\end{pmatrix}.$$
We find $\ww=\begin{pmatrix}0\\1\end{pmatrix}$ and the second solution is $e^{\l t}(\ww+t\vv) = e^{-t}(\begin{pmatrix}0\\1\end{pmatrix}+t\begin{pmatrix}1\\1\end{pmatrix}) = e^{-t}\begin{pmatrix}t\\1+t\end{pmatrix}$.
\com{graphique}
\end{exemple}

\begin{theoreme}
    The linear system is 
    \begin{enumerate}
     \item L-stable if and only if all the eigenvalues of the matrix have non positive real parts, and all the one that are defective have negative real part.
    \item globally asymptotically stable if and only if all the eigenvalues of the matrix
    have negative real parts.
    \end{enumerate}
\end{theoreme}
\begin{proof} 
\quad\\
\begin{enumerate}
\item 
By \prettyref{th:stablecondition}, we just have to show that the conditions on the eigenvalues are equivalent to the fact that the solutions are bounded. The \prettyref{cor:formesolutionlineaire} told us the possible forms of all solutions. The functions $t^je^{\a t}\cos(\b t)$ , $t^je^{\a t}\sin(\b t)$ describe all the possibilities of linear combinations for the coordinates.
Since a linear combination is a finite sum, the solution is bounded if all these possible terms are bounded.

For both, if the eigenvalue $\s=\a+i\b$ is non defective, then $j=0$ only and they are bounded if $\a$ is non positive. If $\s$ is defective, $j>0$ and they are bounded if $\a$ is negative because the exponential is $(t^j)$ for all $j$'s. The sytem is now L-stable.

Alternatively, if there exists an eigenvalue $\s=\a+i\b$ with eigenvector $w=u+iv$ that have a positive real part, then by \prettyref{th:eigensolutions}, there exist a solution $e^{\a t}(cos(\b t)u-\sin(\b t)v)$ and its norm $e^{\a t}\|cos(\b t)u-\sin(\b t)v\|$
is not bounded since $cos(\b t)u-\sin(\b t)v$ does not converge to zero. If $\s$ is defective and is purely imaginary, then \prettyref{th:solutiondegeneree} tell us that there exist a solution $p_w(t)$, polynomial of non null degree, and $p_w(t)\to\infty$ when $t\to\infty$ because of the dominant term. There exists solutions that are unbounded , and the system is not L-stable.

\item 
By \prettyref{th:stablecondition}, we just have to show that the conditions on the eigenvalues are equivalent to the fact that the system is attracting.

But following the considerations of the first part, if all eigenvalues have negative real parts, there is always a $e^{\a t}$ with $\a<0$ term in front of the bases solutions and they all converge to zero. The system is attracting, and globally asymptotically stable.

Alternatively if it is not the case and that $\s=\a+i\b$ has a non negative part, we have a solution $e^{\a t}p_w(t)$ that satisfy $\|e^{\a t}p_w(t)\| \geq \|p_w(t)\|$ which does not go to zero. The system is not attracting and then not globally asymptotically stable.
\end{enumerate}
\end{proof}
\begin{exemple}
     The second condition of the first point seem to point out a very specific situation. Would the theorem work without this condition, meaning that this situation is algebraically impossible? Indeed, we usually do not often see real matrices with a defective and purely imaginary eigenvalue. This is maybe because the following remarks: it can not be in one dimension because eigenvalues are real; it can not be in two dimension because imaginary eigenvalues come by pair of conjugates and then they would be opposed, hence different and non-defective. To have a purely imaginary eigenvalue of multiplicity two we must have its conjugate of multiplicity two too. So we try in four dimension. We search for a matrix with characteristic polynomial $(\l^2+1)^2$. The basis bloc that introduce $i$ eigenvalue is 
     $$J=\begin{pmatrix}0&1\\-1&0\end{pmatrix}.$$
     We could try with diag$(J,J)$, but this gives a non-defective eigenspace. With a little modification on it, we find a matrix that is defective:
     $$\begin{pmatrix}J&I_2\\0&J\end{pmatrix}.$$
     It has the only two independent eigenvectors $(\pm i,1,0,0)^\top$. Hence in this pathological situation, the defective eigensolutions are of the form $e^{it}p(t)=(\cos(t)+i\sin(t))p(t)$, which is neither constant, neither exponential, neither cyclic, neither a polynomial, neither bounded.
\end{exemple}
This is one of the main results of the section, it presents a complete understanding of the stability. It motivates the categorisation of linear systems, taking into account the nature of their stability.
\begin{definition}
    For a sytem $\dotxx=A\xx$, we set the following denomination of the matrix $A$.
    \begin{itemize}
    \item \emph{Stable} : all eigenvalues have negative real part. ($i\R+\R^*_-$)
    \item \emph{unstable} : $-A$ is stable. i.e. all eigenvalues have positive real part. ($i\R+\R^*_+$)
    \item \emph{saddle} : all eigenvalues have a non null real part and they do not all have the same sign.
    \item \emph{hyperbolic} : stable, unstable or saddle; \ie all eigenvalues have non null real part. ($i\R+\R^*$)
    \item \emph{center} : all eigenvalues are purely imaginary. ($i\R^*$)
    \item \emph{elliptic} : all eigenvalues are non null, purely imaginary, and non defective.
    \item \emph{focus} : all eigenvalues are real, negative, and non defective.
    \item \emph{source} : all eigenvalues are real, positive, and non defective.
    \item \emph{degenerated} : there exists a defective eigenvalue.
    \end{itemize}
    More generally, we use these terms to speak about a specific eigenvalue, seeing it as a one-dimensional matrix, meaning that it must respect the generic condition (this does not apply to saddle).
\end{definition}
\begin{remarque}
    The matrix $A$ is unstable (resp. a focus) if $-A$ is stable (resp. a source).
    We point out the similarity between the words "L-stabe" and "stable", which speak respectively about Lyapunov stability and eigenvalues, do not confuse them. A stable linear system is L-stable but not 
\end{remarque}
\com{graphiques et présentations des structures possibles des solutions}
\begin{corollaire}
\quad
\begin{itemize} 
    \item Stable(and focus in particular) and elliptic linear systems are L-stable.
    \item Sources and saddles linear systems are not L-stable, so  neither asymptotically stable.
    \item Stable (and focus in particular) linear systems are globally asymptotically stable.
    \end{itemize} 
\end{corollaire}

All this classification and considerations on the eigenvalues, show us that there exists directions where the trajectories act kind of exponentially, and these are the directions that are linear combinations of some generalised eigenvectors. We present then a separation of the space of solution, which  tell us where the solutions act like in a stable, unstable, or center system.
\begin{definition}
For a linear system $\dotxx=A\xx$ of dimension $n$, we define the stable eigenspace, the unstable eigenspace, and the center eigenspace as follows respectively :
\begin{IEEEeqnarray*}{rCl}
\E^s &=& \Span \{v\in \C^n|\, v\text{ is a generalised eigenvector of the system with eigenvalue }\l \text{ such that }\Re(\l)<0\} \\
\E^u &=& \Span\{v\in \C^n|\, v\text{ is a generalised eigenvector of the system with eigenvalue }\l \text{ such that }\Re(\l)>0\}\\
\E^c &=& \Span\{v\in \C^n|\, v\text{ is a generalised eigenvector of the system with eigenvalue }\l \text{ such that }\Re(\l)=0\}
\end{IEEEeqnarray*}
They are the generalised eigenspaces of the stable, unstable, and center generalised eigenvectors.
In the same idea we want to compare them to the directions where the solutions act exponentially or not:
\begin{IEEEeqnarray*}{rCl}
V^s &=& \{v\in \C^n|\,\exists a,b>0 :\forall t\geq0,\|e^{At}v\| \leq be^{-at}\} 
\\
V^u &=& \{v\in \C^n|\,\exists a,b>0 :\forall t\leq0, \|e^{At}v\| \leq be^{at}\}
\\
V^c &=& \{v\in \C^n|\,\forall a>0,\, e^{At}v e^{-a|t|} \to 0\text{ when }|t|\to\infty\}.
\end{IEEEeqnarray*}
Namely, the spaces where the solutions act exponentially or not.
\end{definition}
\begin{remarque}
We can precise the value of the matrix as an argument when dealing with multiples systems and does not use argument when the context is clear. Note that they are all subspaces and $\E^u(A)=\E^s(-A)$ as well as $V^u(A)=V^s(-A)$. Taking $-A$ instead of $A$ correspond actually to inverse time.
\end{remarque}

\begin{lemme}
    For $\s\in\{s,u,c\}$, we have the equalities $\E^\s=V^\s$, the direct sum $\C^n = \E^s\oplus\E^u\oplus\E^c$, and the $\E^\s$ are invariants in the sense that $\phi(\E^\s\times\R)=\E^\s$, \ie all the solutions stay in the space.
\end{lemme}
\begin{proof}
The invariance of the sets comes from the general form of the eigensolutions presented in \prettyref{th:eigensolutions}. Such a solution looks like $e^{\l t}p(t)$, with an eigenvalue $\l$ and a polynomial $p(t)$ whose coefficients are in the generalised eigenspace of $\l$. It is a linear combination of vectors in the same $\E$ for all $t$. The eigensolutions stay then in the space, as well as the span of them, $\E^\s$.
\\ \\
The direct sum come from the fact that eigenspaces associated to different eigenvalues are independents.The three spaces are just a separation of them in three sums. Namely $\E^\s$ is just the direct sum of the generalised eigenspaces associated to eigenvalues of type $\s$.
\\ \\
Now we prove the equalities. First, we show $\E^\s\subset V^\s$. Each hyperbolic generalised eigensolution associated to a generalised eigenvector $\ww$, can be written $e^{\d t}p(t)$ for a  polynomial $p$ with complex generalised eigenvector coefficients associated to the same eigenvalue $\d=a+ib$ ($a$,$b$ scalars).

$\E^s$: When $a<0$, we can just bound the norm as 
$$\|e^{\d t}p(t)\|=e^{at}\|p(t)\| \leq e^{ta/2} \max_{t\geq0}\||p(t)|e^{ta/2}\| ,$$
The max exists because the continuous expression $p(t)e^{ta/2}\to 0$ when $t\to\infty$ if $a<0$ (by Taylor expansion or iterating Hospital).

$\E^u$: When $a>0$, we directly obtain the $\E^u(A) = \E^s(-A) \subset V^s(-A) = V^u(A)$ by the remark made above.
%When $a>0$ ($\E^u$) and $e^{\d t}p(t)$ is real, we have
%$$\|\Re(e^{\d t}p(t)\| 
%=e^{at}\|p(t))\| 
%\geq e^{at}\min_{\R_+}\|p\|$$
%where the min isn't null otherwise the polynomial would have passed by zero and the trajectory would be identically null.
%If $e^{\d t}p(t)$ isn't real all along, then actually it is never real, otherwise the trajectory would be entirely real by unicity of trajectories. In particular $e^{ibt}p(t)$ is never real and since $\Re(e^{ibt}p(t))>0$. Then
%$$\|\Re(e^{\d t}p(t)\| 
%= e^{at}\|\Re(e^{ibt}p(t)\|
%\geq e^{at/2} \min_{t\geq0}\|e^{at/2}\Re(e^{ibt}p(t))\| .$$
%The min exist since 
%$\|e^{at/2}\Re(e^{ibt}p(t))\|\to \infty$. Indeed $\Re(e^{ibt}p(t))$ couldn't go exponentially to zero since it is constituted of polynimals and trigonometrics terms. The min is non null as we said that $\Re(e^{ibt}p(t))>0$.

$\E^c$: When $a=0$, we have
$$\|e^{-\a |t|}e^{\d t}p(t)\|
= e^{-\a |t|}\|p(t))\| \to0 \text{ when }|t|\to\infty.$$
for any $\a>0$. 

We get then $\ww \in \E^\s$ for the three $\s$. We clearly see that the $V^\s$'s are subspaces, and since $\ww$ was chosen as an arbitrary generalised eigenvector, actually the whole span satisfy $\E^\s\subset V^\s$.

Secondly, we prove the inverse inclusion.
%\com{first try} The $V^\s$'s contain independents subspaces $\E^\s$ whose dimensions sum up to the dimension of the all space. As a result, if one of them was strictly bigger, we would have a dependency between them for example a linear combination $0 = \b_s v_s + \b_u v_u + \b_c v_c$ with the complex $v_\s$ not all null. Then $0 = \|\b_s v_s + \b_u v_u + \b_c v_c\|$.
%\com{second try}
Each vector $v$ of $V^\s$ can be written as a unique sum of vectors of the three $\E^\s$: $v = v_s + v_u + v_c$. We define a new norm on $\C^n$ using this unique writing. 
$$N(v)= \big\|(\|v_s\|,\|v_u\|,\|v_c\|)\big\|_1 = \|v_s\|+\|v_u\|+\|v_c\|.$$ Indeed, absolute homogeneity and positive definiteness are clear, while subadditivity come from the unique decomposition: for vectors $v$ and $w$, we have $v_s+w_s + v_u+w_u + v_c+w_c=v+u$. This implies necessarily that $(u+v)_\s=v_\s + w_\s$, and then
$$N(v+w) = \|v_s+w_s\|+\|v_u+w_u\|+\|v_c+w_c\| \leq \|v_s\|+\|w_s\|+\|v_u\|+\|w_u\|+\|v_c\|+ \|w_c\|= N(v)+N(w).$$
Now, $N$ is equivalent to any other norm, \ie we have a constant $c$ such that $c^{-1}\|v\|\leq N(v)\leq c\|v\|$.
If $\s=s$, $N(e^{At}v)\to0$ exponentially when $t\to\infty$. By looking at the norm $N$, we see that it must be the same for the directions $v_s$, $v_u$ and $v_c$. From the general form of eigensolutions we see that unstable and center eigensolutions does not even goes to zero: 
$$\|e^{\d t}p(t)\|
=\|e^{\Re(\d)t}p(t)\|
\geq e^{\Re(\d)t}\min\|p\|\geq \min\|p\|>0$$
Where the min exist and is positive because the polynomial is not trivial. Then actually $v_u=v_c=0$, and $v = v_s\in \E^s$, so $V^s\subset \E^s$. As before we directly conclude that $$V^u(A)=V^s(-A)\subset \E^s(-A)=\E^u(A).$$ Same kind of argument for a $v=v_s+v_u+v_c\in V^c$: 
$$e^{-\a|t|}\|e^{\d t}p(t)\|
=e^{-\a|t|}e^{at}\|p(t)\|
\geq e^{(at/|t|-\a)|t|}\min\|p\|\geq \min\|p\|>0$$
for a small $0<\a<|a|$ and alternatively for positive time if $\d$ unstable ($a>0$), or for negative time if $\d$ stable ($a<0$). We get that $v$ is in $\E^c$.
We have proven the both inclusions and obtain the result.
\end{proof}
Now that we understand well the complex eigenspaces of solutions, we can easily deal with the real case:
\begin{corollaire}
    We have the equalities $\Re(\E^\s)=\Re(V^\s)$ and $\dim(\Re(\E^\s))=\dim(\E^\s)$, the direct sum
    $\R^n = \Re(\E^s)\oplus\Re(\E^u)\oplus\Re(\E^c)$, and the sets $\E^s$ are invariant:
    $\phi(\Re(\E^\s)\times\R)=\Re(\E^\s)$.
\end{corollaire}
\begin{proof}
The equality $\Re(\E^\s)=\Re(V^\s)$ is trivial. We have the simple sum
$$\R^n 
=\Re(\C^n) 
=\Re(\E^s\oplus\E^u\oplus\E^c)
= \Re(\E^s)+\Re(\E^u)+\Re(\E^c).$$
This is actually a direct sum because the conjugate invariance of the $\E^\s$'s implies that $\Re(\E^\s)\subset\E^\s$. This last inclusion implies that $\dim(\Re(\E^\s))\leq\dim(\E^\s)$ and this must be a equality because the dimensions must sum up to $n$. The invariance of the spaces along solutions come from the fact that a real solution stay real in the time and $\E^\s$ is invariant.
\end{proof}
\com{Graphic of a $\R^3$ hyperbolic system with the stability eigenspaces (a plane and a line)}
\begin{remarque}
By taking the real part of the spaces (just like the imaginary part) and summing them, we can construct all the real space and the induced real solutions. However, note that taking the real part of a basis will not a priory give a real basis. Indeed the complex basis can use complex coefficient to build all real solutions while a real basis can only use real coefficients and the real part of a product is not just the product of the real parts. 
\end{remarque}
We see that the eigenspaces seem to act independently from each other and the general comportment of solutions in them, seem to be somewhat uniform. A change of an eigenvalue can change totally the stability of a system; the flow $e^{At}x$ is continuous on $A$, but how does the uniformity continuity change along $t$ and $x$? We can show that actually, in the hyperbolic case, the dimensions of the eigenspaces do not change when we chose a second matrix near enough the first one, in some sense the general stability structure depends continuously on $A$:

\begin{theoreme}
For $\s=s,u$, the map $A\mapsto \text{dim }\E^\s(A)$ is locally continuous on hyperbolic matrices, and so locally constant.
\end{theoreme}
\begin{proof}
We treat only the stable eigenspace since $\text{dim }\E^u(A) =\text{dim }\E^s(-A)$. It is a classical but non trivial analysis result that the roots of a polynomial are continuous, so we prove that if we can find a neighbourhood where the roots stay stable, the structure of the generalised eigenspaces stay the same. Since no eigenvalue is in the imaginary line, we can surround all the stable eigenvalues by a positively oriented simple closed rectifiable curve $\g$, that does not cross any eigenvalue, the imaginary line neither. Since the characteristic polynomial $p_A$ vanishes exactly on the eigenvalues, and is holomorphic, we have by the residues that 
$$\frac{1}{2\pi i}\int_{\g} \frac{p_A'}{p_A}
= \text{dim }\E^s(A) .$$
Indeed, if $p_A$ has a root $\l$ of order $k$, the residue of $p'_A/p_a$ in $\l$ is just $k$. Since $p_A$ is a polynomial, the order of a zero is just the its multiplicity, and it correspond then to the dimension of the generalised eigenspace of this eigenvalue. The space $\E^\s$ is a direct sum of the generalised eigenspaces of stable eigenvalues, whose dimensions sum up to the dimension of the stable eigenspace itself. Since the coefficients of $p_A$ are polynomials of the coefficients of $A$, they change continuously with $A$. We can suppose that $A$ varies in a small neighbourhood so $p_A$ still do not vanish in the curves $\g$ (that does not change with $A$). As a result, the formula is well defined around $A$ and is continuous with respect to $p_A$ and hence to $A$.
\end{proof}
 As we see, the stability structure of the system is the same for matrices in a neighbourhood. But more generally, if two matrices $A$ and $B$ are far from each other, but still have the same dimensions of eigenspaces, do them have similarities? We introduce the following tool of comparison, that match all the trajectories from both system in a continuous way. 
 \begin{definition}
 Two flows $\phi$ and $\psi$ on the space $\R^n$, are said \emph{topologically conjugate}, if there exist a homeomorphism $h:\R^n\to\R^n$ such that 
 $h\circ \phi_t = \psi_t\circ h$, i.e. the image of a trajectory is a trajectory. 
 \end{definition}
 \begin{lemme} \label{lem:ps}
If a matrix is stable with maximum real part of eigenvalues $-a<0$, there exists a scalar product $\ps{.}{.}_a$ such that $\ps{x}{Ax}_a<-a\|x\|_a^2$.
\end{lemme}
\begin{proof}
Linear algebra result, see \com{ref} for a proof.
\end{proof}
\begin{theoreme}[Hartman and Grobman]
Two stable near systems on the same space are topologically conjugate.
\end{theoreme}
\begin{proof}
To find a homeomorphism $h$, we first search one on a subset of the space that content points that are part of all non trivial solutions and then we will extend it along the trajectories. The idea is to find a kind of sphere that trajectories cross only one time. The \prettyref{lem:ps} exhibit a special scalar product that seem to be related to the problem. We consider two matrices $A$, $B$ and their associated linear systems. We set the spheres $S_A$ and $S_B$ associated to the specials products $\ps{}{}_A$, $\ps{}{}_B$ of the matrices $A$ and $B$. We show that trajectories cross them exactly one time by showing that the norm chosen is monotone along the solutions:
$$\ddt \|\xx\|_A
= \ddt \ps{\xx}{\xx}_A
= 2\ps{\xx}{\dotxx}_A
= 2\ps{\xx}{A\xx}_A
\leq -a\|\xx\|^2_A$$
For a scalar $a>0$. So before applying Grönwal and conclude the exponential convergence to zero, we see that this convergence is strictly monotone definite in the sense of the norm $\|.\|_A$. Thus, $\xx$ only passes in $S_A$ one time. The same hold for $B$ by symetry of the situation. We can set now an homeomorphism 
\begin{IEEEeqnarray*}{rCl}
h_0:S_A &\to& S_B \\
\xx &\mapsto& \xx/\|\xx\|_B
\end{IEEEeqnarray*}
It is well defined because the image is obviously in $S_B$, $\xx$ is never 0 and the denominator never vanishes. The continuity come from continuity of the norm. By symmetry of the situation its inverse should be $g_0:\xx \mapsto \xx/\|\xx\|_A$. Indeed for a $\xx\in S_B$,
$$ (h_0\circ g_0)(\xx) = h_0(\frac{\xx}{\|\xx\|_A}) 
= \frac{\frac{\xx}{\|\xx\|_A}}{\|\frac{\xx}{\|\xx\|_A}\|_B}
= \frac{\xx}{\|\xx\|_B} =\xx$$
so $h_0\circ g_0 = \text{Id}_{S_B}$ and by the symetry of the situation, $g_0\circ h_0 = \text{Id}_{S_A}$. The function $h_0$ is a continuous bijection, an homeomorphism.

We construct now the final $h$. For this we want to take a initial point, find the intersection of its trajectory with the sphere, go to the other sphere, and return go back in the trajectory as much as we came in the first sphere. For this we need to know for any point $\xx$, how many time $\tau(x)$ does it take to go to the sphere, \ie $e^{A\tau(x)}x\in S_A$. The value $\tau(x)$ is uniquely defined by the choice of the sphere. We set $h$ by
$$h(\xx) 
= \begin{cases} 
    0 \text{ if } \xx=0 \\
    e^{-B\tau(\xx)}h_0(e^{A\tau(\xx)}\xx) \text{ else}
  \end{cases}
$$
We prove first that it has indeed the relation $h(e^{At}\xx)=e^{Bt}h(\xx)$. We note that $\tau(e^{At}x)= \tau(x)-t$ because $$\phi(e^{At}x,\tau(x)-t)
= e^{A(\tau(x)-t)}e^{At}x
= e^{A\tau(\xx)}x \in S_A.$$ When $\xx$ is not the origin we compute,
\begin{IEEEeqnarray*}{rCl}
h(e^{At}\xx) 
&=& e^{-B\tau(e^{At}\xx)} h_0(e^{A\tau(e^{At}\xx)}e^{At}\xx) \\
&=& e^{-B(\tau(\xx)-t)} h_0(e^{A(\tau(\xx)-t)}e^{At}\xx)\\
&=& e^{Bt}e^{-B\tau(\xx)} h_0(e^{A\tau(\xx)}\xx) \\
&=& e^{Bt} h(\xx)
\end{IEEEeqnarray*}
If we write it $h=h_{AB}$, we show that it has an inverse $g=h_{BA}$. By looking at the definition of $h(\xx)$, we see that 
$$e^{B\tau(\xx)}h(\xx) = h_0(e^{A\tau(\xx)}\xx) \in S_B$$
so $\tau(h(\xx))=\tau(\xx)$. We compute that for all $\xx\in S_A$
\begin{IEEEeqnarray*}{rCl}
g\circ h(\xx)
&=& e^{-A\tau(h(\xx))}g_0(e^{B\tau(h(\xx))}h(\xx)) \\
&=& e^{-A\tau(\xx)}g_0(e^{B\tau(\xx)}h(\xx)) \\
&=& e^{-A\tau(\xx)}g_0(e^{B\tau(\xx)}e^{-B\tau(\xx)}h_0(e^{A\tau(\xx)}\xx))\\
&=& e^{-A\tau(\xx)}g_0(h_0(e^{A\tau(\xx)}\xx)) \\
&=& e^{-A\tau(\xx)}e^{A\tau(\xx)}\xx \\
&=& x
\end{IEEEeqnarray*}
and $g\circ h= \text{Id}_{S_A}$. By symmetry of the situation, $h\circ g= \text{Id}_{S_B}$ too and $h$ is a bijection with inverse $g$.

It remind to show the continuity of $h$, which is the same for $g$. Apart of the origin, everything is continuous if we can prove it for $\tau$. We put the case orgine away and will do it separately and for now we can use the implicit function theorem on $(\xx,t)\mapsto\|e^{At}\xx\|_A-1$, which is continuous. The partial derivative with respect to t and in a point $(\xx_0,t_0)$ is 
$$\partial_t(\|e^{At_0}\xx_0\|_A-1)
=2\ps{e^{At_0}\xx_0}{Ae^{At_0}\xx_0}
<-a\|e^{At_0}\| < 0$$
and is inversible because monotone. We must have then a continuous map $\tilde{\tau}$ in a neighbourhood of $\xx_0$ such that $\|e^{A\tilde{\tau}(\xx)}\xx\|_A-1=0$. By unicity of $\tau$ (and $\tilde{\tau}$ actually), we know that in this neighbourhood, $\tau=\tilde{\tau}$ is continuous. Since $\xx_0$ was chosen arbitrarily, $\tau$ is locally continuous and so continuous. For $\xx\neq0$, $h(\xx)$ is a composition of continuous function and is itself continuous.

For the continuity in $\xx=0$, we take a sequence $(\xx_k)_k\to0$  and we have necessarily that $\tau(\xx_k)$ goes to $-\infty$. If not, then we must have a bounded sub-sequence $\xx_{k_j}$ since $\tau$ is necessarily negative in the unit ball. The result is that we can suppose this sub-sequence converging to a time $\tau^*$, and get by the continuity of the flow
$$S_A \ni \phi(\xx_{k_j},\tau(\xx_{k_j}))
\xrightarrow[j\to\infty]{} \phi(0,\tau^*)=0,$$ which is a contradiction with the closure of the sphere. We use $\tau(\xx_k)\to-\infty$ to get 
$$\|h(\xx_k)\|_B 
= \|e^{-B\tau(\xx_k)}h_0(e^{A\tau(\xx_k)}\xx_k)\|_B 
\leq \|e^{-B\tau(\xx_k)}\|_B\to 0$$
because $B$ is stable. Finally, $h$ is continuous in the origin and everywhere and it is a homeomorphism.
\end{proof}
This theorem give us a total correspondence between two flows, and show how the eigenvalues are decisive in the structure of the stability. We can actually show that the unstable and saddle cases have a similar result, built on the theorem. We present a partial proof to see what is going on. 
\begin{corollaire}
    Two hyperbolic linear systems with stable eigenspace of same dimension, are topologically conjugates.
\end{corollaire}
\begin{proof}
    Let $A$ and $B$ be unstable matrices with induced linear flows $\phi$ and $\psi$. The linear systems of matrices $-A$ and $-B$ are stable, they are topologically conjugate by a homeomorphism $h$. Now simply $$h\circ\phi_t(\xx) 
    = h(e^{At}\xx) 
    = h(e^{-A(-t)}\xx) 
    = e^{-B(-t)}h(\xx)
    = e^{Bt}h(\xx)
    = \phi_t\circ h(\xx)
    ,$$
    so the same $h$ works for the unstable case. We do not prove here the saddle case, but we present a sketch of proof:
    When the two matrices are saddle, but with stable eigenspaces of the same dimension $k$ and $m$, we can decompose the space like
    $$\R^n
    = \E^s(A)+\E^u(A) 
    =\E^s(B)+\E^u(B)$$
    Finite and same dimensional spaces are homeomorphic by a bijective linear map so we can use the theorem in each eigenspace independently as it was $\R^d$ with $d$ the dimension of the stable or unstable eigenspace. We would get two homeomorphism $h_s:\E^s(A)\to\E^s(B)$ and $h_u:\E^u(A)\to\E^u(B)$. To define a $h$ in the all space we can use the projections $\pi_s$ and $\pi_u$ of the decomposition $\R^n
    = \E^s(A)+\E^u(A)$  and we would get something like $h = h_s\circ\pi_s + h_u\circ\pi_u$ that will give us a homeomorphism.
\end{proof}
We have studied the stability of linear dynamical system under multiples point of view. Linear system can be though to be a very specific situation. It is indeed. However, we present here a first link to the non-linear case, whose stability appear to be determined in some way by a linear associated system. This motivate the study of linear system and their possible applications to non-linear dynamics.
\begin{theoreme}[Linearization, Hartman–Grobman] \label{th:linearisation}
    Consider a system $\dotxx=\FF(\xx)$ such that $\FF$ is derivable, and $\xx^*$ a fixed point. If the linearized system $\dotxx=D\FF(\xx^*)\xx$ is stable, then $\xx^*$ is asymptotically stable for the non linear system.
\end{theoreme}
\begin{proof}
If the fixed point is not the origin, we set $\yy=\xx-\xx^*$ which satisfy $\dotyy=\dotxx=\FF(\xx)=\FF(\yy+\xx^*)$. Its flow is just a translation of the flow of the initial problem, is therefore isometric to it, and has the corresponding fixed point in the origin. Since $\FF$ is $C^2$, we can write its Taylor expression around the origin:
$$\FF(\xx)= \FF(0)+D\FF(\xx)\xx + h(\xx)\xx =: A\xx + h(\xx)\xx.$$
for some function $h(\xx)\to0$ when $\xx\to0$. We use now the special scalar product $\ps{.}{.}_a$ related to the stable matrix $A$ of \prettyref{lem:ps} and the metric induced by it, which is equivalent to the euclidean metric. We get
\begin{IEEEeqnarray*}{rCl}
\ddt \ps{\xx}{\xx}_a
&=& 2\ps{\xx}{\dotxx}_a
= 2\ps{\xx}{\FF(\xx)}_a 
= 2\ps{\xx}{A\xx + h(\xx)\xx}_a
= 2\ps{\xx}{A\xx}_a + 2\ps{\xx}{h(\xx)\xx}_a \\
&\leq& -a\|\xx\|_a^2 + 2\|\xx\|^2_ah(\xx)\\
&=& (2h(x)- a)\|\xx\|_a^2. \\
&\leq&  -c\|\xx\|_a^2
\end{IEEEeqnarray*}
We can take $\|\xx\|$ small enough such that $2h(x)-a$ is negative, let us say for any $a>\e>0$, $\|\xx\|<\d(\e)$ implies $\a=2h(x)-a<-\e$ Then by Grönwall, $\ddt\|\xx\|^2_a<-\a\|\xx\|^2_a$ implies 
$$\|\xx\|^2_a < \|\xx(0)\|_a^2e^{-\a t}\to0$$
exponentially when $t\to\infty$ and the system is attracting since it happen in a neighbourhood $B_{\|.\|_a}(0,\d(\e))$ of the origin. Note that taking a neighbourhood small enough, we can attain a exponential convergence of any rate $\a<a$. The L-stability come just from the fact that the convergence of the attractivity is direct and does not go out its chosen neighbourhood, namely a sphere for the metric $\|.\|_a$. This give us the asymptotic stability.
\end{proof}
The idea behind this proof is that we have defined a quantity, $\|\xx\|_a$, and that we have showed that this quantity was always decreasing along the trajectories, like a potential energy. This kind of function is called a Lyapunov function, we will define it in the next chapter and see that it has nice general applications.
    \chapter{deuxième}
Lorem ipsum dolor sit amet, consectetur adipiscing elit. Sed non risus. Suspendisse lectus tortor, dignissim sit amet, adipiscing nec, ultricies sed, dolor. Cras elementum ultrices diam. Maecenas ligula massa, varius a, semper congue, euismod non, mi. Proin porttitor, orci nec nonummy molestie, enim est eleifend mi, non fermentum diam nisl sit amet erat. Duis semper. Duis arcu massa, scelerisque vitae, consequat in, pretium a, enim. Pellentesque congue. Ut in risus volutpat libero pharetra tempor. Cras vestibulum bibendum augue. Praesent egestas leo in pede. Praesent blandit odio eu enim. Pellentesque sed dui ut augue blandit sodales. Vestibulum ante ipsum primis in faucibus orci luctus et ultrices posuere cubilia Curae; Aliquam nibh. Mauris ac mauris sed pede pellentesque fermentum. Maecenas adipiscing ante non diam sodales hendrerit.

Ut velit mauris, egestas sed, gravida nec, ornare ut, mi. Aenean ut orci vel massa suscipit pulvinar. Nulla sollicitudin. Fusce varius, ligula non tempus aliquam, nunc turpis ullamcorper nibh, in tempus sapien eros vitae ligula. Pellentesque rhoncus nunc et augue. Integer id felis. Curabitur aliquet pellentesque diam. Integer quis metus vitae elit lobortis egestas. Lorem ipsum dolor sit amet, consectetuer adipiscing elit. Morbi vel erat non mauris convallis vehicula. Nulla et sapien. Integer tortor tellus, aliquam faucibus, convallis id, congue eu, quam. Mauris ullamcorper felis vitae erat. Proin feugiat, augue non elementum posuere, metus purus iaculis lectus, et tristique ligula justo vitae magna.

Aliquam convallis sollicitudin purus. Praesent aliquam, enim at fermentum mollis, ligula massa adipiscing nisl, ac euismod nibh nisl eu lectus. Fusce vulputate sem at sapien. Vivamus leo. Aliquam euismod libero eu enim. Nulla nec felis sed leo placerat imperdiet. Aenean suscipit nulla in justo. Suspendisse cursus rutrum augue. Nulla tincidunt tincidunt mi. Curabitur iaculis, lorem vel rhoncus faucibus, felis magna fermentum augue, et ultricies lacus lorem varius purus. Curabitur eu amet.
    \chapter{troisième}
Lorem ipsum dolor sit amet, consectetur adipiscing elit. Sed non risus. Suspendisse lectus tortor, dignissim sit amet, adipiscing nec, ultricies sed, dolor. Cras elementum ultrices diam. Maecenas ligula massa, varius a, semper congue, euismod non, mi. Proin porttitor, orci nec nonummy molestie, enim est eleifend mi, non fermentum diam nisl sit amet erat. Duis semper. Duis arcu massa, scelerisque vitae, consequat in, pretium a, enim. Pellentesque congue. Ut in risus volutpat libero pharetra tempor. Cras vestibulum bibendum augue. Praesent egestas leo in pede. Praesent blandit odio eu enim. Pellentesque sed dui ut augue blandit sodales. Vestibulum ante ipsum primis in faucibus orci luctus et ultrices posuere cubilia Curae; Aliquam nibh. Mauris ac mauris sed pede pellentesque fermentum. Maecenas adipiscing ante non diam sodales hendrerit.

Ut velit mauris, egestas sed, gravida nec, ornare ut, mi. Aenean ut orci vel massa suscipit pulvinar. Nulla sollicitudin. Fusce varius, ligula non tempus aliquam, nunc turpis ullamcorper nibh, in tempus sapien eros vitae ligula. Pellentesque rhoncus nunc et augue. Integer id felis. Curabitur aliquet pellentesque diam. Integer quis metus vitae elit lobortis egestas. Lorem ipsum dolor sit amet, consectetuer adipiscing elit. Morbi vel erat non mauris convallis vehicula. Nulla et sapien. Integer tortor tellus, aliquam faucibus, convallis id, congue eu, quam. Mauris ullamcorper felis vitae erat. Proin feugiat, augue non elementum posuere, metus purus iaculis lectus, et tristique ligula justo vitae magna.

Aliquam convallis sollicitudin purus. Praesent aliquam, enim at fermentum mollis, ligula massa adipiscing nisl, ac euismod nibh nisl eu lectus. Fusce vulputate sem at sapien. Vivamus leo. Aliquam euismod libero eu enim. Nulla nec felis sed leo placerat imperdiet. Aenean suscipit nulla in justo. Suspendisse cursus rutrum augue. Nulla tincidunt tincidunt mi. Curabitur iaculis, lorem vel rhoncus faucibus, felis magna fermentum augue, et ultricies lacus lorem varius purus. Curabitur eu amet. 
    
    \nocite{*}
    \printbibliography
\end{document}
